% physics_hitchin_langlands.tex
%
% Physics note: Higgs bundles, the Hitchin system, and the geometric Langlands
% programme from the quantum vertex chiral group perspective.
%
% Not included in main.tex; standalone research note.
%
% Raeez Lorgat, April 2026

\documentclass[11pt]{article}
\usepackage[T1]{fontenc}
\usepackage[utf8]{inputenc}
\usepackage{lmodern}
\frenchspacing
\usepackage[cmintegrals,cmbraces,noamssymbols]{newtxmath}
\usepackage{ebgaramond}
\let\Bbbk\undefined
\let\openbox\undefined
\usepackage{amsmath,amssymb,amsthm}
\usepackage{mathtools}
\usepackage{enumitem}
\usepackage{tikz-cd}
\usepackage[unicode,pdfencoding=auto,psdextra]{hyperref}
\usepackage{geometry}
\geometry{top=1.2in, bottom=1.2in, left=1.25in, right=1.25in}

% Theorem environments
\newtheorem{theorem}{Theorem}[section]
\newtheorem{proposition}[theorem]{Proposition}
\newtheorem{lemma}[theorem]{Lemma}
\newtheorem{corollary}[theorem]{Corollary}
\newtheorem{conjecture}[theorem]{Conjecture}
\theoremstyle{definition}
\newtheorem{definition}[theorem]{Definition}
\newtheorem{example}[theorem]{Example}
\newtheorem{construction}[theorem]{Construction}
\theoremstyle{remark}
\newtheorem{remark}[theorem]{Remark}
\newtheorem{warning}[theorem]{Warning}

% Macros (subset of main.tex)
\newcommand{\Higgs}{\mathrm{Higgs}}
\newcommand{\Hit}{\mathrm{Hit}}
\newcommand{\Bun}{\mathrm{Bun}}
\newcommand{\Loc}{\mathrm{Loc}}
\newcommand{\Rep}{\mathrm{Rep}}
\newcommand{\Coh}{\mathrm{Coh}}
\newcommand{\Fuk}{\mathrm{Fuk}}
\newcommand{\Mod}{\mathrm{Mod}}
\newcommand{\Hom}{\mathrm{Hom}}
\newcommand{\End}{\mathrm{End}}
\newcommand{\Sym}{\mathrm{Sym}}
\newcommand{\Tr}{\mathrm{Tr}}
\newcommand{\Ran}{\mathrm{Ran}}
\newcommand{\Spec}{\mathrm{Spec}}
\newcommand{\id}{\mathrm{id}}
\newcommand{\frakg}{\mathfrak{g}}
\newcommand{\frakh}{\mathfrak{h}}
\newcommand{\frakn}{\mathfrak{n}}
\newcommand{\HH}{\mathrm{HH}}
\newcommand{\HC}{\mathrm{HC}}
\newcommand{\CY}{\mathrm{CY}}
\newcommand{\DT}{\mathrm{DT}}
\newcommand{\GV}{\mathrm{GV}}
\newcommand{\BPS}{\mathrm{BPS}}
\newcommand{\cA}{\mathcal{A}}
\newcommand{\cB}{\mathcal{B}}
\newcommand{\cC}{\mathcal{C}}
\newcommand{\cD}{\mathcal{D}}
\newcommand{\cM}{\mathcal{M}}
\newcommand{\cN}{\mathcal{N}}
\newcommand{\cO}{\mathcal{O}}
\newcommand{\cZ}{\mathcal{Z}}
\newcommand{\cH}{\mathcal{H}}
\newcommand{\Prym}{\mathrm{Prym}}
\newcommand{\Jac}{\mathrm{Jac}}
\newcommand{\Eone}{E_1}
\newcommand{\Etwo}{E_2}
\DeclareMathOperator{\rk}{rk}
\DeclareMathOperator{\mult}{mult}

\begin{document}

\begin{center}
{\LARGE\scshape Higgs Bundles, the Hitchin System, and\\
Geometric Langlands from Quantum Vertex Chiral Groups}

\bigskip

{\large Raeez Lorgat}

\medskip

{\itshape Research note --- April 2026}
\end{center}

\bigskip

\tableofcontents

\bigskip

\section*{Introduction}

This note develops the connection between the Hitchin moduli space of Higgs bundles on a curve and the quantum vertex chiral group framework of Volume~III.  The central claim is that the Hitchin system for a reductive group $G$ on a curve $C$ of genus $g$ is a Calabi--Yau geometry whose quantum vertex chiral group $\mathcal{G}(C, G)$ encodes the geometric Langlands correspondence: Koszul duality of quantum vertex chiral groups exchanges $G$ and $G^\vee$.

The note is organised as follows.  Section~\ref{sec:hitchin-moduli} recalls the Hitchin moduli space as a non-compact hyperk\"ahler manifold and establishes its Calabi--Yau structure.  Section~\ref{sec:hitchin-fibration} develops the completely integrable structure.  Section~\ref{sec:qvcg-hitchin} constructs the quantum vertex chiral group $\mathcal{G}(C, G)$ and analyses its root datum.  Section~\ref{sec:langlands-duality} formulates the Langlands duality conjecture at the level of chiral groups.  Section~\ref{sec:6d-compactification} derives the construction from the 6d $(2,0)$ theory.


% ==========================================================================
\section{The Hitchin moduli space}
\label{sec:hitchin-moduli}
% ==========================================================================

\subsection{$G$-Higgs bundles}

Fix a smooth projective curve $C$ of genus $g \geq 2$ and a complex reductive group $G$ with Lie algebra $\frakg$.

\begin{definition}[$G$-Higgs bundle]
\label{def:higgs-bundle}
A \emph{$G$-Higgs bundle} on $C$ is a pair $(E, \varphi)$ where:
\begin{enumerate}[label=(\roman*)]
  \item $E$ is a principal $G$-bundle on $C$;
  \item $\varphi \in H^0(C, \mathrm{ad}(E) \otimes K_C)$ is a section of the adjoint bundle twisted by the canonical bundle --- the \emph{Higgs field}.
\end{enumerate}
\end{definition}

The Higgs field $\varphi$ is a $(1,0)$-form valued in $\frakg$; it can be viewed as a deformation of the Dolbeault operator $\bar\partial_E$ to the flat ``Higgs connection'' $\bar\partial_E + \varphi$.  The condition that $(E, \varphi)$ be semistable is the vanishing of the moment map for the gauge group action (the Hitchin equation).

\begin{definition}[The Hitchin moduli space]
\label{def:hitchin-moduli}
The \emph{Hitchin moduli space} is the moduli space of semistable $G$-Higgs bundles on $C$:
\[
  \cM_H(C, G) = \bigl\{ (E, \varphi) \mid E \text{ a principal } G\text{-bundle}, \; \varphi \in H^0(C, \mathrm{ad}(E) \otimes K_C), \; (E, \varphi) \text{ semistable} \bigr\} / \sim
\]
where $\sim$ denotes $S$-equivalence.
\end{definition}

\subsection{Dimension and hyperk\"ahler structure}

The Hitchin moduli space has a rich differential-geometric structure descending from Hitchin's self-duality equations.

\begin{proposition}[Hitchin]
\label{prop:hitchin-hk}
The smooth locus of $\cM_H(C, G)$ is a hyperk\"ahler manifold of real dimension
\[
  \dim_{\mathbb{R}} \cM_H(C, G) = 4 \dim(G)(g - 1).
\]
\end{proposition}

\begin{proof}[Sketch]
The tangent space at a smooth point $(E, \varphi)$ is the first hypercohomology $\mathbb{H}^1$ of the deformation complex
\[
  \mathrm{ad}(E) \xrightarrow{[\varphi, -]} \mathrm{ad}(E) \otimes K_C.
\]
By Riemann--Roch, $\dim_\mathbb{C} \mathbb{H}^1 = 2\dim(G)(g-1)$, hence the real dimension is $4\dim(G)(g-1)$.  The hyperk\"ahler structure comes from interpreting $\cM_H$ as Hitchin's moduli space of solutions to the self-duality equations
\[
  F_A + [\varphi, \varphi^*] = 0, \qquad \bar\partial_A \varphi = 0,
\]
which is a hyperk\"ahler quotient of the infinite-dimensional affine space of connections and Higgs fields by the gauge group.
\end{proof}

The three complex structures $(I, J, K)$ give three descriptions:
\begin{enumerate}[label=(\roman*)]
  \item \textbf{Complex structure $I$} (Dolbeault): $\cM_H$ parametrises Higgs bundles $(E, \varphi)$;
  \item \textbf{Complex structure $J$} (de Rham): $\cM_H \cong \cM_{\mathrm{dR}}(C, G)$, the moduli of flat $G$-connections (via the non-abelian Hodge correspondence);
  \item \textbf{Complex structure $K$} (Betti): $\cM_H \cong \cM_B(C, G) = \Hom(\pi_1(C), G)/G$, the character variety.
\end{enumerate}

\subsection{Calabi--Yau structure}

\begin{proposition}
\label{prop:hitchin-cy}
The Hitchin moduli space $\cM_H(C, G)$ is a (non-compact) Calabi--Yau manifold.  More precisely:
\begin{enumerate}[label=(\roman*)]
  \item The holomorphic symplectic form $\omega_I$ (in complex structure $I$) is
  \[
    \omega_I = \int_C \Tr(\delta\varphi \wedge \delta A)
  \]
  where $\delta A$ and $\delta\varphi$ are tangent vectors.  This is nowhere-degenerate.
  \item The canonical bundle is trivial: $K_{\cM_H} \cong \cO_{\cM_H}$.  The trivialising section is the top exterior power of $\omega_I$.
  \item The holonomy is contained in $\mathrm{Sp}(n, \mathbb{C}) \subset \mathrm{SU}(2n)$ where $n = \dim(G)(g-1)$, hence in particular in $\mathrm{SU}(2n)$.
\end{enumerate}
\end{proposition}

\begin{remark}[CY dimension]
\label{rem:cy-dim-hitchin}
As a complex manifold (in complex structure $I$), $\cM_H$ has complex dimension $2\dim(G)(g-1)$.  It is Calabi--Yau of this complex dimension (holomorphic symplectic, with holonomy $\mathrm{Sp}(n) \subset \mathrm{SU}(2n)$ where $n = \dim(G)(g-1)$).  For $G = \mathrm{SL}_2$ and $g = 2$: $\dim_\mathbb{C} = 6$, so $\cM_H$ is a (non-compact) holomorphic symplectic sixfold with $\mathrm{Sp}(3)$ holonomy.  The Hitchin fibration gives a Lagrangian fibration of ``CY3 type'': both the base and the generic fibre have complex dimension~3.  This is the critical dimension for the quantum vertex chiral group construction (see Remark~\ref{rem:cy3-realisation}).
\end{remark}

\begin{warning}
\label{warn:non-compact}
$\cM_H$ is non-compact.  This is essential: the Hitchin fibration (Section~\ref{sec:hitchin-fibration}) gives a $\mathbb{C}^*$-action $\varphi \mapsto t\varphi$ whose fixed locus is $\Bun_G(C)$ (the zero section).  The non-compactness means that the CY category $D^b(\Coh(\cM_H))$ requires appropriate support conditions (e.g.\ ind-coherent sheaves, or restriction to the nilpotent cone).
\end{warning}


% ==========================================================================
\section{The Hitchin fibration}
\label{sec:hitchin-fibration}
% ==========================================================================

\subsection{The Hitchin base and the characteristic polynomial}

The Hitchin fibration is the map that sends a Higgs bundle to the ``characteristic polynomial'' of its Higgs field.

\begin{definition}[Hitchin base]
\label{def:hitchin-base}
Let $p_1, \ldots, p_r$ be homogeneous generators of the invariant ring $\mathbb{C}[\frakg]^G$, of degrees $d_1, \ldots, d_r$ (where $r = \rk(G)$).  The \emph{Hitchin base} is the affine space
\[
  \cA_H = \bigoplus_{i=1}^r H^0(C, K_C^{d_i}).
\]
Its dimension is $\dim \cA_H = \sum_{i=1}^r (2d_i - 1)(g - 1) = \dim(G)(g-1)$, which is half the (complex) dimension of $\cM_H$.
\end{definition}

\begin{definition}[Hitchin fibration]
\label{def:hitchin-fibration}
The \emph{Hitchin fibration} is
\[
  \pi \colon \cM_H(C, G) \longrightarrow \cA_H, \qquad (E, \varphi) \longmapsto (p_1(\varphi), \ldots, p_r(\varphi)).
\]
\end{definition}

\begin{example}[$G = \mathrm{GL}_n$]
\label{ex:hitchin-gln}
For $G = \mathrm{GL}_n$, the invariants are $p_i = \Tr(\wedge^i)$, so $\pi(E, \varphi) = (\Tr(\varphi), \Tr(\varphi^2), \ldots, \det(\varphi))$.  The Hitchin base is $\cA_H = \bigoplus_{i=1}^n H^0(C, K_C^i)$.  For $G = \mathrm{SL}_2$: $\cA_H = H^0(C, K_C^2)$, the space of quadratic differentials, of dimension $3g - 3$.
\end{example}

\subsection{Complete integrability}

\begin{theorem}[Hitchin]
\label{thm:hitchin-integrable}
The Hitchin fibration $\pi \colon \cM_H(C, G) \to \cA_H$ is a completely integrable Hamiltonian system with respect to the holomorphic symplectic form $\omega_I$.  That is:
\begin{enumerate}[label=(\roman*)]
  \item $\dim \cA_H = \frac{1}{2} \dim_\mathbb{C} \cM_H$;
  \item the fibres of $\pi$ are Lagrangian with respect to $\omega_I$;
  \item the functions $p_i(\varphi)$ Poisson-commute.
\end{enumerate}
\end{theorem}

\subsection{The spectral curve and Prym varieties}

\begin{definition}[Spectral curve]
\label{def:spectral-curve}
For a point $a = (a_1, \ldots, a_r) \in \cA_H$, the \emph{spectral curve} $\widetilde{C}_a$ is the curve in the total space of $K_C$ defined by the characteristic equation
\[
  \det(\eta \cdot \id - \varphi) = \eta^n + a_1 \eta^{n-1} + \cdots + a_n = 0
\]
(for $G = \mathrm{GL}_n$; more generally, $\widetilde{C}_a$ is the cameral cover determined by $a$).  The projection $p \colon \widetilde{C}_a \to C$ is a degree-$n$ branched cover.
\end{definition}

\begin{proposition}
\label{prop:hitchin-fiber}
For a generic point $a \in \cA_H$ (where $\widetilde{C}_a$ is smooth):
\begin{enumerate}[label=(\roman*)]
  \item The Hitchin fibre $\pi^{-1}(a)$ is isomorphic to the Prym variety $\Prym(\widetilde{C}_a / C)$ (or a torsor thereof), an abelian variety of dimension $\dim(G)(g-1)$;
  \item For $G = \mathrm{GL}_n$: $\pi^{-1}(a) \cong \Jac(\widetilde{C}_a)$, the Jacobian of the spectral curve;
  \item For $G = \mathrm{SL}_n$: $\pi^{-1}(a) \cong \Prym(\widetilde{C}_a / C)$, the kernel of the norm map $\mathrm{Nm} \colon \Jac(\widetilde{C}_a) \to \Jac(C)$.
\end{enumerate}
\end{proposition}

The Hitchin system is thus a family of abelian varieties (generically) fibred over an affine base.  The singular fibres --- where the spectral curve degenerates --- carry the interesting topology.

\begin{remark}[The discriminant locus]
\label{rem:discriminant}
The locus $\Delta \subset \cA_H$ where the spectral curve is singular is the \emph{discriminant divisor}.  Over the complement $\cA_H \setminus \Delta$, the Hitchin fibration is a smooth torus fibration.  The monodromy of this torus fibration around $\Delta$ is a central ingredient in the Strominger--Yau--Zaslow picture of mirror symmetry (Section~\ref{sec:syz}).
\end{remark}


% ==========================================================================
\section{The quantum vertex chiral group $\mathcal{G}(C, G)$}
\label{sec:qvcg-hitchin}
% ==========================================================================

We now apply the general construction of Volume~III to the Calabi--Yau geometry $\cM_H(C, G)$.

\subsection{When $\cM_H$ is CY3}

The complex dimension of $\cM_H(C, G)$ is $2\dim(G)(g - 1)$.  For the quantum vertex chiral group to arise from a CY3, we need
\[
  2\dim(G)(g - 1) = 3, \quad \text{i.e.} \quad \dim(G)(g-1) = \frac{3}{2}.
\]
This is satisfied by $G = \mathrm{SL}_2$ (so $\dim(G) = 3$) and $g = 2$ (giving $\frac{3}{2} \cdot 2 = 3$), or more properly by the quotient construction in Remark~\ref{rem:cy3-realisation} below.  In general, $\cM_H$ is a higher-dimensional CY manifold, and we work with the CY-$d$ generalisation of the quantum vertex chiral group from Section~\ref{sec:higher-cy-generalisation}.

\begin{remark}[CY3 realisation]
\label{rem:cy3-realisation}
There are two natural ways to obtain a CY3 from the Hitchin system for arbitrary $G$:
\begin{enumerate}[label=(\alph*)]
  \item \textbf{Local CY3}: Consider the cotangent bundle $T^*\Bun_G(C)$ (or a suitable open subset).  In the Dolbeault description, $\cM_H$ is a deformation of $T^*\Bun_G$.  The fibre of $T^*\Bun_G$ over a smooth point has complex dimension $\dim(G)(g-1)$, which need not be 3.  However, one can consider a \emph{local model}: for each spectral curve $\widetilde{C}_a$, the local CY3 geometry is $T^*\widetilde{C}_a$ (complex dimension 2, not 3) or a suitable 3-fold thickening.
  \item \textbf{Dimensional reduction}: Compactify the 6d $(2,0)$ theory on $C$ to get a 4d $\cN = 4$ theory with Coulomb branch $\cM_H$.  Further compactify on $S^1$ to get 3d $\cN = 4$, whose target is $\cM_H$.  The CY3 structure appears in the \emph{effective} geometry of the 2d sigma model (after a second $S^1$ compactification), as developed in Section~\ref{sec:6d-compactification}.
\end{enumerate}
For the quantum vertex chiral group construction, we use interpretation (b): the partition function of the 2d sigma model into $\cM_H$ is a character of $\mathcal{G}(C, G)$, and the CY3 structure arises from the effective geometry rather than the literal complex dimension of $\cM_H$.
\end{remark}

\subsection{The root datum $\mathcal{R}(C, G)$}

\begin{construction}[Root datum of $\mathcal{G}(C, G)$]
\label{constr:root-datum-hitchin}
The generalised root datum $\mathcal{R}(C, G)$ of the quantum vertex chiral group associated to $\Higgs(C, G)$ consists of:

\begin{enumerate}[label=(\roman*)]
  \item \textbf{Lattice.}  The lattice is
  \[
    \Lambda(C, G) = K_0(\Coh(\cM_H(C, G))) \cong \Lambda_{\mathrm{cochar}}(G) \oplus H_1(C, \mathbb{Z})
  \]
  where $\Lambda_{\mathrm{cochar}}(G)$ is the cocharacter lattice of $G$ and $H_1(C, \mathbb{Z}) \cong \mathbb{Z}^{2g}$ encodes the topology of $C$.  The bilinear form on $\Lambda(C, G)$ combines the Killing form on $\Lambda_{\mathrm{cochar}}$ with the intersection form on $H_1(C)$.

  \item \textbf{Real roots.}  The real simple roots are the roots of $G$ itself:
  \[
    \Delta^{\mathrm{re}} = \Phi(G) = \text{the root system of } G.
  \]
  These are encoded in the Cartan matrix $A_{ij} = \langle \alpha_i, \alpha_j^\vee \rangle$ and have $(\alpha, \alpha) > 0$.  They generate the ``finite'' part of the Weyl group.

  \item \textbf{Imaginary roots from the Hitchin base.}  The imaginary roots arise from the Hitchin base $\cA_H$: for each degree $d_i$ ($i = 1, \ldots, r = \rk(G)$) and each element of $H^0(C, K_C^{d_i})$, there is an imaginary root.  More precisely, the imaginary root multiplicities are controlled by a \emph{spectral generating function}
  \[
    \sum_{\alpha \in \Delta^{\mathrm{im}}} \mult(\alpha) \, q^\alpha = \sum_{i=1}^r \chi(C, K_C^{d_i}) \cdot \delta_i
  \]
  where $\chi(C, K_C^{d_i}) = (2d_i - 1)(g - 1)$ and $\delta_i$ are the imaginary simple roots corresponding to the $i$-th Casimir invariant.

  \item \textbf{BPS multiplicities.}  The root multiplicities are DT/BPS invariants of $\cM_H$:
  \[
    \mult(\alpha) = \Omega(\alpha; C, G) = \text{BPS index at charge } \alpha.
  \]
  For the real roots, $\mult(\alpha) = 1$.  For the imaginary roots, the multiplicities encode the ``number of BPS states'' in the Hitchin fibre --- equivalently, the cohomology of the Hitchin fibre at charge $\alpha$.

  \item \textbf{Weyl group.}  The Weyl group is
  \[
    W(C, G) = W(G) \ltimes \Lambda_{\mathrm{cochar}}(G)^{2g}
  \]
  where $W(G)$ is the finite Weyl group of $G$ and $\Lambda_{\mathrm{cochar}}(G)^{2g}$ acts by lattice translations indexed by $H_1(C, \mathbb{Z})$.  This is an affine Weyl group of genus-$g$ type.

  \item \textbf{Weyl vector.}  The Weyl vector $\rho(C, G) \in \frakh^* \otimes \mathbb{Q}$ satisfies
  \[
    \langle \rho, \alpha_i^\vee \rangle = 1 \quad \text{for all simple roots } \alpha_i \in \Delta^{\mathrm{re}}.
  \]
  This is the usual Weyl vector of $G$, shifted by a genus-dependent correction from the Hitchin base.
\end{enumerate}
\end{construction}

\begin{remark}[Comparison with $K3 \times E$]
\label{rem:comparison-k3e}
For $G = \mathrm{SL}_2$, $g = 2$, the root datum $\mathcal{R}(C, G)$ should be compared with the root datum of the $K3 \times E$ tower (the $K3 \times E$ chapter of Vol.~III):
\begin{itemize}
  \item The real roots $\Phi(\mathrm{SL}_2) = \{\pm\alpha\}$ play the role of $\{\delta_1, \delta_2, \delta_3\}$ in the $K3 \times E$ story.
  \item The Hitchin base $\cA_H = H^0(C, K_C^2) \cong \mathbb{C}^3$ has the same dimension as the positive cone in $\Lambda^{2,1}_{II}$.
  \item The imaginary root multiplicities are governed by DT invariants of $\cM_H(\mathrm{SL}_2)$ rather than by the K3 elliptic genus $\phi_{0,1}$ --- but in both cases, the multiplicities are Fourier coefficients of an automorphic form.
\end{itemize}
\end{remark}

\subsection{The denominator identity as an automorphic form for $G^\vee$}

The denominator identity of the BKM superalgebra $\frakg_{\cM_H}$ associated to the root datum $\mathcal{R}(C, G)$ should be an automorphic form for the \emph{Langlands dual group} $G^\vee$.

\begin{conjecture}[Automorphic denominator identity]
\label{conj:automorphic-denom}
The Weyl--Kac--Borcherds denominator identity of the generalised BKM superalgebra $\frakg_{\mathcal{R}(C,G)}$ takes the form
\[
  \Phi_G(z) = e^{-2\pi i (\rho, z)} \prod_{\alpha \in \Delta_+} (1 - e^{-2\pi i (\alpha, z)})^{\mult(\alpha)}
\]
and $\Phi_G$ is an automorphic form for a suitable arithmetic subgroup $\Gamma \subset G^\vee(\mathbb{Z})$.
\end{conjecture}

The evidence:
\begin{enumerate}[label=(\alph*)]
  \item For $K3 \times E$ with the $\mathrm{SL}_2$ Hitchin system at $g = 2$: the denominator identity is $\Delta_5$, the Igusa cusp form, which is a Siegel modular form for $\mathrm{Sp}_4(\mathbb{Z})$.  Here $\mathrm{Sp}_4 \cong \mathrm{SO}_+(3,2) \cong \mathrm{SO}_+(\Lambda^{3,2})$, and the Langlands dual of $\mathrm{SL}_2$ is $\mathrm{PGL}_2$ --- the group whose Siegel-type dual is $\mathrm{Sp}_4$.
  \item Gritsenko--Nikulin's general theory: reflective automorphic forms on $\mathrm{O}(2, n)$ are denominator identities of generalised BKM algebras, and the arithmetic group is determined by the lattice, hence by $G^\vee$.
  \item The Langlands dual structure is visible in the exponents: the degrees $d_1, \ldots, d_r$ of the Casimir invariants of $G$ are the same as those of $G^\vee$ (since $G$ and $G^\vee$ have the same Weyl group), ensuring $\dim \cA_H(C, G) = \dim \cA_H(C, G^\vee)$.
\end{enumerate}


\subsection{Higher-CY generalisation}
\label{sec:higher-cy-generalisation}

For $\dim_\mathbb{C}(\cM_H) = 2\dim(G)(g-1) > 3$, the moduli space is CY of dimension $d > 3$.  The quantum vertex chiral group generalises as follows.

\begin{construction}[CY-$d$ quantum vertex chiral group]
\label{constr:cyd-qvcg}
For a CY category $\cC$ of dimension $d$, the quantum vertex chiral group $\mathcal{G}(\cC)$ is defined by the same pipeline as in the CY3 case (root datum from BPS spectrum, automorphic correction via shadow obstruction tower), but:
\begin{enumerate}[label=(\roman*)]
  \item The $\mathbb{S}^d$-framing replaces the $\mathbb{S}^3$-framing, giving $E_d \to E_2$ by Dunn additivity (rather than $E_3 \to E_2$).
  \item The shadow obstruction tower has $d$-dependent depth: the modular characteristic $\kappa(\mathcal{G}(\cC))$ involves $\lambda_g$ on $\overline{\cM}_{g,1}$ weighted by $\hbar^{2g-2}$.
  \item The BKM superalgebra may be replaced by a higher BKM structure (a vertex algebra with $d$-shifted Lie bracket).
\end{enumerate}
\end{construction}


% ==========================================================================
\section{Langlands duality as BKM Koszul duality}
\label{sec:langlands-duality}
% ==========================================================================

This section contains the central conjecture: the geometric Langlands correspondence is realised by Koszul duality of quantum vertex chiral groups.

\subsection{The classical geometric Langlands correspondence}

We recall the main players.

\begin{enumerate}[label=(\roman*)]
  \item The \textbf{stack of $G$-bundles} $\Bun_G(C)$ and its derived category $D^b(\Coh(\Bun_G))$.
  \item The \textbf{stack of $G^\vee$-local systems} $\Loc_{G^\vee}(C)$ parametrising flat $G^\vee$-connections (or equivalently $G^\vee$-representations of $\pi_1(C)$) and its category of quasi-coherent sheaves $\mathrm{QCoh}(\Loc_{G^\vee})$.
  \item The \textbf{geometric Langlands conjecture} (Beilinson--Drinfeld, Arinkin--Gaitsgory, Fargues--Scholze in the $p$-adic setting, proved by Gaitsgory et al.\ 2024):
  \[
    D\text{-}\Mod(\Bun_G(C)) \simeq \mathrm{IndCoh}_{\cN}(\Loc_{G^\vee}(C))
  \]
  as dg-categories, where the left side is $D$-modules on $\Bun_G$ and the right side is ind-coherent sheaves on $\Loc_{G^\vee}$ with nilpotent singular support.
\end{enumerate}

The non-abelian Hodge correspondence identifies $\Loc_{G^\vee}(C)$ with $\cM_H(C, G^\vee)$ (in complex structure $J$).  Thus the Langlands correspondence relates:
\begin{itemize}
  \item The ``$A$-side'': $D$-modules on $\Bun_G \subset \cM_H(C, G)$ (the zero section in complex structure $I$);
  \item The ``$B$-side'': coherent sheaves on $\cM_H(C, G^\vee)$ (in complex structure $J$, alias flat $G^\vee$-connections).
\end{itemize}

\subsection{SYZ mirror symmetry for the Hitchin system}
\label{sec:syz}

The Hitchin system provides a geometric realisation of the Langlands correspondence via mirror symmetry.

\begin{proposition}[Hausel--Thaddeus, Donagi--Pantev]
\label{prop:syz-hitchin}
The Hitchin fibrations for $G$ and $G^\vee$ are \emph{SYZ mirror duals}:
\[
  \begin{tikzcd}
    \cM_H(C, G) \arrow[dr, "\pi_G"'] & & \cM_H(C, G^\vee) \arrow[dl, "\pi_{G^\vee}"] \\
    & \cA_H &
  \end{tikzcd}
\]
Over the regular locus $\cA_H^{\mathrm{reg}}$:
\begin{enumerate}[label=(\roman*)]
  \item The fibres $\pi_G^{-1}(a)$ and $\pi_{G^\vee}^{-1}(a)$ are dual abelian varieties;
  \item T-duality (fibrewise Fourier--Mukai transform) along the Hitchin fibres exchanges $G$ and $G^\vee$;
  \item This T-duality is the geometric incarnation of Langlands duality.
\end{enumerate}
\end{proposition}

The key point: the Hitchin base $\cA_H$ is the \emph{same} for $G$ and $G^\vee$ (since the Casimir degrees are the same), but the fibres are dual tori.  The exchange of the torus and its dual is classical mirror symmetry / T-duality.

\subsection{The Koszul duality conjecture}

\begin{conjecture}[Langlands duality as BKM Koszul duality]
\label{conj:langlands-koszul}
The quantum vertex chiral groups $\mathcal{G}(C, G)$ and $\mathcal{G}(C, G^\vee)$ are Koszul dual:
\[
  \boxed{\mathcal{G}(C, G)^! = \mathcal{G}(C, G^\vee).}
\]
More precisely:
\begin{enumerate}[label=(\roman*)]
  \item The Verdier dual of the bar complex of $\mathcal{G}(C, G)$ (as an $E_2$-factorization coalgebra) is identified with the bar complex of $\mathcal{G}(C, G^\vee)$:
  % AP25: The Koszul dual is obtained via Verdier duality D_Ran(B(A))
  % = B(A^!), NOT via the cobar functor Omega(B(A)) = A (which
  % recovers the ORIGINAL algebra by bar-cobar inversion).
  \[
    D_{\mathrm{Ran}}(B_{E_2}(\mathcal{G}(C, G))) \simeq B_{E_2}(\mathcal{G}(C, G^\vee)).
  \]
  \item At the level of root data, Koszul duality exchanges:
  \begin{itemize}
    \item the root lattice $\Lambda_{\mathrm{root}}(G)$ with the coroot lattice $\Lambda_{\mathrm{coroot}}(G) = \Lambda_{\mathrm{root}}(G^\vee)$;
    \item the real roots $\Phi(G)$ with the coroots $\Phi^\vee(G) = \Phi(G^\vee)$;
    \item the imaginary root multiplicities $\mult(\alpha; G)$ with $\mult(\alpha^\vee; G^\vee)$.
  \end{itemize}
  \item The denominator identities are related by the functional equation of the automorphic form under the Langlands dual group.
  \item The modular characteristics satisfy
  \[
    \kappa(\mathcal{G}(C, G)) + \kappa(\mathcal{G}(C, G^\vee)) = \kappa_{\mathrm{crit}}
  \]
  where $\kappa_{\mathrm{crit}}$ is the critical value --- the analogue of $c + c' = 26$ for Virasoro, or $K + K' = K_{\mathrm{crit}}$ for the appropriate W-algebra.
\end{enumerate}
\end{conjecture}

\begin{remark}[Relation to the complementarity principle]
\label{rem:complementarity-langlands}
Conjecture~\ref{conj:langlands-koszul}(iv) is an instance of the complementarity principle from Volume~I: the quantum vertex chiral group of $G$ and its Koszul dual (the quantum vertex chiral group of $G^\vee$) have modular characteristics summing to a critical value.  For Virasoro ($G = \mathrm{SL}_2$), the critical value is $c + c' = 26$; for Kac--Moody $\widehat{\frakg}_k$, it is $k + k' = k_{\mathrm{crit}}$ where $k_{\mathrm{crit}} = -h^\vee$.  The Hitchin/Langlands generalisation replaces the level by the full modular characteristic $\kappa$.
\end{remark}

\subsection{Evidence and consistency checks}

\begin{enumerate}[label=\textbf{(\arabic*)}]
  \item \textbf{Abelian case $G = T$ (a torus).}  For a torus, $\cM_H(C, T) = T^*\Bun_T(C) = T^*(T^{2g})$ and $G^\vee = T^\vee$ is the dual torus.  The quantum vertex chiral groups are abelian (no real roots), and Koszul duality is Fourier--Mukai duality on the torus --- the classical abelian Langlands correspondence.

  \item \textbf{$G = \mathrm{SL}_2$, $g = 2$.}  As noted in Remark~\ref{rem:comparison-k3e}, this case overlaps with the $K3 \times E$ story.  The denominator identity $\Delta_5$ is self-dual under $\mathrm{Sp}_4(\mathbb{Z})$, consistent with $\mathrm{SL}_2^\vee = \mathrm{PGL}_2$ (and $\mathrm{SL}_2 \cong \mathrm{PGL}_2$ at the level of Lie algebras).

  \item \textbf{Simply-laced $G$.}  If $G$ is simply laced, then $G \cong G^\vee$, so $\mathcal{G}(C, G)^! \cong \mathcal{G}(C, G)$: the quantum vertex chiral group is \emph{self-dual}.  This is the analogue of the self-duality of the Hitchin system for simply-laced groups (the spectral curves and their duals are isomorphic).

  \item \textbf{Non-simply-laced: $B_n / C_n$ duality.}  $\mathrm{SO}(2n+1)^\vee = \mathrm{Sp}(2n)$.  The Koszul duality $\mathcal{G}(C, \mathrm{SO}(2n+1))^! = \mathcal{G}(C, \mathrm{Sp}(2n))$ exchanges short and long roots, i.e.\ exchanges the root system of $B_n$ with the coroot system $C_n$.

  \item \textbf{Quantum geometric Langlands.}  Gaitsgory's quantum geometric Langlands (QGL) conjecture identifies twisted $D$-modules at level $\kappa$ on $\Bun_G$ with twisted $D$-modules at the ``complementary'' level $\kappa^\vee$ on $\Bun_{G^\vee}$.  In the quantum vertex chiral group language, QGL becomes the statement that the $\Etwo$-chiral representation categories are equivalent:
  \[
    \Rep^{\Etwo}_\kappa(\mathcal{G}(C, G)) \simeq \Rep^{\Etwo}_{\kappa^\vee}(\mathcal{G}(C, G^\vee))
  \]
  with $\kappa + \kappa^\vee = \kappa_{\mathrm{crit}}$.
\end{enumerate}


% ==========================================================================
\section{The 6d $(2,0)$ perspective}
\label{sec:6d-compactification}
% ==========================================================================

The deepest physical justification for the structures above comes from the 6d $(2,0)$ superconformal field theory.

\subsection{The 6d $(2,0)$ theory of type $G$}

For each simply-laced Lie algebra $\frakg$ of type ADE, there exists a 6d superconformal field theory $\mathcal{T}^{6d}_\frakg$ --- the $(2,0)$ theory --- defined on a 6-manifold $M_6$.  Its defining properties:
\begin{enumerate}[label=(\roman*)]
  \item It has $(2,0)$ supersymmetry: 16 supercharges, with R-symmetry $\mathrm{Sp}(4)_R \cong \mathrm{SO}(5)_R$.
  \item The tensor branch moduli space is $(\mathbb{R}^5)^r / W(\frakg)$ where $r = \rk(\frakg)$.
  \item On its Coulomb branch, it is a free theory of $r$ self-dual 2-form fields $B^i_+$ ($i = 1, \ldots, r$).
  \item It has no Lagrangian description (for $\frakg \neq \mathfrak{su}(1)$).
\end{enumerate}

For non-simply-laced $G$, the $(2,0)$ theory is replaced by a ``class $\mathcal{S}$'' construction involving outer automorphism twists.

\subsection{Compactification on $C$: 4d $\cN = 2$ theories of class $\mathcal{S}$}

Compactifying $\mathcal{T}^{6d}_\frakg$ on $C$ (with a partial topological twist along $C$) gives a 4d $\cN = 2$ theory $\mathcal{T}^{4d}[C, \frakg]$ --- a ``theory of class $\mathcal{S}$'' in the sense of Gaiotto.

\begin{proposition}
\label{prop:class-s}
The Coulomb branch of $\mathcal{T}^{4d}[C, \frakg]$ is the Hitchin base $\cA_H$, and the total space of the Seiberg--Witten fibration is the Hitchin moduli space:
\[
  \mathcal{B}_{\mathrm{Coulomb}}(\mathcal{T}^{4d}[C, \frakg]) = \cA_H, \qquad \mathcal{M}_{\mathrm{SW}} = \cM_H(C, G).
\]
The Seiberg--Witten curve at a point $a \in \cA_H$ is the spectral curve $\widetilde{C}_a$.
\end{proposition}

\subsection{Further compactification: 3d $\cN = 4$ and 2d sigma models}

\begin{construction}[The compactification cascade]
\label{constr:cascade}
\phantom{x}
\begin{enumerate}[label=\textbf{Step \arabic*.}]
  \item \textbf{$6d \to 4d$}: Compactify $\mathcal{T}^{6d}_\frakg$ on $C$ (topological twist).  Result: 4d $\cN = 2$ theory $\mathcal{T}^{4d}[C, G]$ with target $\cM_H(C, G)$.

  \item \textbf{$4d \to 3d$}: Compactify further on $S^1_R$ of radius $R$.  Result: 3d $\cN = 4$ sigma model with target $\cM_H(C, G)$.  In the limit $R \to 0$: the 3d mirror is the sigma model into $\cM_H(C, G^\vee)$.  \emph{This is 3d mirror symmetry = Langlands duality at the level of hyperk\"ahler targets.}

  \item \textbf{$3d \to 2d$}: Compactify on a second $S^1_r$ of radius $r$.  Result: 2d $\cN = (4,4)$ sigma model with target $\cM_H(C, G)$.  Taking the $A$-twist gives the $A$-model; taking the $B$-twist gives the $B$-model.

  \item \textbf{$2d \to$ partition function}: The partition function of the 2d sigma model on a genus-$h$ Riemann surface $\Sigma_h$ is
  \[
    Z(\Sigma_h; \cM_H(C, G)) = \Tr_{\cH} \, q^{L_0} \bar{q}^{\bar{L}_0}
  \]
  where $\cH$ is the Hilbert space of the 2d theory.
\end{enumerate}
\end{construction}

\subsection{The partition function as a character of $\mathcal{G}(C, G)$}

\begin{conjecture}[Partition function = character]
\label{conj:partition-function}
The partition function of the 2d sigma model into $\cM_H(C, G)$ on a torus $\Sigma_1 = T^2$ is a character of the quantum vertex chiral group $\mathcal{G}(C, G)$:
\[
  Z(T^2; \cM_H(C, G)) = \mathrm{ch}(\mathcal{G}(C, G)) = \sum_{\alpha} \mult(\alpha) \, q^\alpha.
\]
More precisely:
\begin{enumerate}[label=(\roman*)]
  \item The graded dimension of the BKM superalgebra $\frakg_{\mathcal{R}(C,G)}$ equals the elliptic genus of $\cM_H(C, G)$.
  \item The denominator identity of $\frakg_{\mathcal{R}(C,G)}$ equals the partition function of the $B$-twisted sigma model.
  \item The genus-$h$ partition function $Z(\Sigma_h; \cM_H)$ is a genus-$h$ correlation function of $\mathcal{G}(C, G)$.
\end{enumerate}
\end{conjecture}

The chain of identifications:
\[
  \begin{array}{rcl}
    \text{6d $(2,0)$ on $C \times T^2$} & \longleftrightarrow & \text{2d sigma model into } \cM_H(C, G) \\[4pt]
    \text{partition function } Z(T^2; \cM_H) & \longleftrightarrow & \text{character of } \mathcal{G}(C, G) \\[4pt]
    \text{denominator identity of BKM} & \longleftrightarrow & \text{automorphic form for } G^\vee \\[4pt]
    \text{Langlands duality } G \leftrightarrow G^\vee & \longleftrightarrow & \text{Koszul duality } \mathcal{G}(C,G)^! = \mathcal{G}(C, G^\vee)
  \end{array}
\]

\begin{remark}[Comparison with $K3 \times E$]
\label{rem:6d-k3e}
The $K3 \times E$ tower arises from a different compactification: the IIA string on $K3 \times E$ (or equivalently, the heterotic string on $T^5$, via string duality).  The 2d worldsheet theory of the IIA string has a partition function controlled by $\Delta_5$.  In the Hitchin language, this corresponds to $G = \mathrm{SL}_2$, $g = 2$, where the sigma model target $\cM_H(C, \mathrm{SL}_2)$ is a (non-compact) CY3.  The BKM superalgebra $\frakg_{\Delta_5}$ IS the quantum vertex chiral group $\mathcal{G}(C, \mathrm{SL}_2)$ for this case.
\end{remark}


% ==========================================================================
\section{Detailed structures}
\label{sec:detailed-structures}
% ==========================================================================

\subsection{The spectral curve as a BPS generating function}

The spectral curve $\widetilde{C}_a$ over a point $a \in \cA_H$ encodes the BPS spectrum at that point of the Coulomb branch.  In the quantum vertex chiral group language:

\begin{proposition}
\label{prop:spectral-bps}
The spectral curve $\widetilde{C}_a$ determines the imaginary root multiplicities of $\mathcal{G}(C, G)$ at ``charge'' $a$:
\begin{enumerate}[label=(\roman*)]
  \item Over the generic locus: $\widetilde{C}_a$ is smooth, and the Hitchin fibre $\Prym(\widetilde{C}_a/C)$ is an abelian variety whose cohomology determines $\mult(\alpha)$ for $\alpha$ in the fibre direction.
  \item At the discriminant locus: $\widetilde{C}_a$ degenerates, new BPS states appear (wall-crossing), and the root multiplicities jump.
  \item The total multiplicity function $\alpha \mapsto \mult(\alpha)$ is piecewise constant on chambers of $\cA_H \setminus \Delta$, with wall-crossing formula governing the jumps across $\Delta$.
\end{enumerate}
\end{proposition}

\begin{remark}[Wall-crossing and gauge equivalence]
\label{rem:wall-crossing}
The Kontsevich--Soibelman wall-crossing formula, which governs the jumps in BPS multiplicities across walls of marginal stability, should be interpreted in the quantum vertex chiral group framework as \emph{gauge equivalence of MC elements}: the shadow obstruction tower $\Theta_{\cA}^{\leq r}$ at different points of the Coulomb branch $\cA_H$ are related by gauge transformations, and the wall-crossing automorphism is the transition function.  This is consistent with Open Question~6 of the Vision document.
\end{remark}

\subsection{The Feigin--Frenkel center and the Hitchin base}

There is a deep connection between the Hitchin base and the center of the completed enveloping algebra at the critical level.

\begin{proposition}[Feigin--Frenkel]
\label{prop:ff-center}
The center of the completed enveloping algebra $\widehat{U}(\widehat{\frakg})_{\mathrm{crit}}$ at the critical level $k = -h^\vee$ is isomorphic to the algebra of functions on the space of $G^\vee$-opers on the formal disc:
\[
  \mathfrak{z}(\widehat{\frakg})_{\mathrm{crit}} \cong \mathrm{Fun}(\mathrm{Op}_{G^\vee}(\mathcal{D})).
\]
Globalising to the curve $C$: the space of $G^\vee$-opers $\mathrm{Op}_{G^\vee}(C)$ embeds into $\Loc_{G^\vee}(C) \cong \cM_H(C, G^\vee)$ (in complex structure $J$), and the Hitchin map for $G^\vee$ restricted to opers recovers the ``oper base'' inside $\cA_H$.
\end{proposition}

In the quantum vertex chiral group language:
\begin{itemize}
  \item The Feigin--Frenkel center $\mathfrak{z}(\widehat{\frakg})_{\mathrm{crit}}$ is the \emph{center} of the quantum vertex chiral group $\mathcal{G}(C, G)$ at the critical level $\kappa = \kappa_{\mathrm{crit}}$.
  \item The space of opers is the ``classical limit'' of $\mathcal{G}(C, G^\vee)$ --- the base of the Hitchin fibration for $G^\vee$.
  \item The Feigin--Frenkel isomorphism is a shadow of the Koszul duality $\mathcal{G}(C, G)^! = \mathcal{G}(C, G^\vee)$: at the critical level, the center of one side becomes the spectrum of the other.
\end{itemize}

\subsection{The $\Etwo$-structure and the Kapustin--Witten twist}

The $\Etwo$-enhancement of the quantum vertex chiral group has a physical origin in the Kapustin--Witten topological twist.

\begin{construction}[Kapustin--Witten twist]
\label{constr:kw-twist}
The 4d $\cN = 4$ super Yang--Mills theory (obtained by compactifying the $(2,0)$ theory on $T^2$) admits a family of topological twists parametrised by $t \in \mathbb{CP}^1$ (the Kapustin--Witten parameter).  The two extreme cases are:
\begin{enumerate}[label=(\roman*)]
  \item $t = 0$: the $B$-twist.  BPS equations: $F_A = 0$ (flat connections).  The category of branes is $D$-$\Mod(\Bun_G)$.
  \item $t = \infty$: the $A$-twist.  BPS equations: Hitchin equations.  The category of branes is $\Fuk(\cM_H(C, G))$.
  \item $t = 1$: the ``GL twist'' (geometric Langlands twist).  This is the twist relevant for the geometric Langlands correspondence.
\end{enumerate}
At $t = 1$, the braided monoidal structure on the category of line operators is the $\Etwo$-structure of $\mathcal{G}(C, G)$.
\end{construction}

\begin{remark}
\label{rem:e2-from-4d}
The $\Etwo$-structure arises because line operators in a 4d topological theory can be braided (moved around each other in the transverse $\mathbb{R}^2$), giving a braided monoidal category.  This is exactly the $\Etwo$-chiral structure of the quantum vertex chiral group: the braiding comes from the $\Etwo$-operad acting on the transverse directions.
\end{remark}


% ==========================================================================
\section{The Hitchin system for specific groups}
\label{sec:specific-groups}
% ==========================================================================

\subsection{$G = \mathrm{SL}_2$: the simplest case}

For $G = \mathrm{SL}_2$:
\begin{itemize}
  \item $\dim(G) = 3$, $r = \rk(G) = 1$, Casimir degree $d_1 = 2$.
  \item $\cA_H = H^0(C, K_C^2)$, dimension $3(g-1)$.
  \item $\cM_H(C, \mathrm{SL}_2)$ has $\dim_\mathbb{C} = 6(g-1)$.
  \item For $g = 2$: $\dim_\mathbb{C} = 6$ (holomorphic symplectic sixfold).  The Hitchin fibration is a Lagrangian fibration of CY3 type: base $\cA_H \cong \mathbb{C}^3$, generic fibre $\Prym \cong$ abelian 3-fold.
  \item Root datum: $\Delta^{\mathrm{re}} = \{\pm \alpha\}$ (the root system of $\mathrm{SL}_2$), imaginary roots from $H^0(C, K_C^2)$.
  \item $G^\vee = \mathrm{PGL}_2$ (at the group level; at the Lie algebra level, $\mathfrak{sl}_2 \cong \mathfrak{pgl}_2$, so the root datum is self-dual).
\end{itemize}

The spectral curve for $\mathrm{SL}_2$ at $a \in H^0(C, K_C^2)$ is
\[
  \widetilde{C}_a \colon \eta^2 + a = 0 \quad \subset \text{Tot}(K_C),
\]
a double cover of $C$ branched at the $4(g-1)$ zeros of $a$.  For $g = 2$: 4 branch points.  By Riemann--Hurwitz, $2g(\widetilde{C}_a) - 2 = 2(2 \cdot 2 - 2) + 4 = 8$, so $g(\widetilde{C}_a) = 5$.  The Prym variety $\Prym(\widetilde{C}_a/C)$ has dimension $g(\widetilde{C}_a) - g(C) = 5 - 2 = 3$.

\subsection{$G = \mathrm{SL}_n$: spectral curves and Jacobians}

For $G = \mathrm{SL}_n$:
\begin{itemize}
  \item Casimir degrees $d_i = i+1$ for $i = 1, \ldots, n-1$.
  \item $\cA_H = \bigoplus_{i=2}^n H^0(C, K_C^i)$, dimension $(n^2-1)(g-1)$.
  \item Spectral curve: $\eta^n + a_2\eta^{n-2} + \cdots + a_n = 0$ in Tot$(K_C)$ (trace-free condition eliminates the $\eta^{n-1}$ term).
  \item Generic Hitchin fibre: $\Prym(\widetilde{C}_a / C)$, an abelian variety of dimension $(n^2-1)(g-1)$.
  \item Langlands dual: $G^\vee = \mathrm{PGL}_n$.  Root lattice of $\mathrm{SL}_n$ is dual to coroot lattice of $\mathrm{PGL}_n$ --- this exchanges the Hitchin fibres as dual abelian varieties (Hausel--Thaddeus for $n = 2$; Donagi--Pantev in general).
\end{itemize}

\subsection{$G = E_8$: the maximally exotic case}

For $G = E_8$ (simply laced, hence self-dual):
\begin{itemize}
  \item $\dim(G) = 248$, $r = 8$, Casimir degrees $2, 8, 12, 14, 18, 20, 24, 30$.
  \item $\dim_\mathbb{C} \cM_H = 496(g-1)$: a CY manifold of very high dimension.
  \item The quantum vertex chiral group $\mathcal{G}(C, E_8)$ is self-dual: $\mathcal{G}(C, E_8)^! = \mathcal{G}(C, E_8)$.
  \item The denominator identity should be an automorphic form for $E_8(\mathbb{Z})$ of weight determined by the genus $g$.
  \item For $g = 2$: the automorphic form lives on a domain of very high rank.  The construction of this form (and verification of the BKM structure) is an open problem.
\end{itemize}


% ==========================================================================
\section{The dictionary}
\label{sec:dictionary}
% ==========================================================================

We summarise the correspondence between the various perspectives.

\medskip

\begin{center}
\renewcommand{\arraystretch}{1.4}
\begin{tabular}{@{}p{3.5cm}p{3.5cm}p{4.2cm}@{}}
\hline
\textbf{Hitchin / physics} & \textbf{Geometric Langlands} & \textbf{Quantum vertex chiral group} \\
\hline
$G$-Higgs bundle $(E, \varphi)$ & $G$-bundle on $C$ & Generator of $\mathcal{G}(C,G)$ \\
Hitchin base $\cA_H$ & Opers on $C$ & Imaginary root space \\
Spectral curve $\widetilde{C}_a$ & Oper at point $a$ & BPS charge lattice at $a$ \\
Hitchin fibre $\Prym(\widetilde{C}_a/C)$ & Hecke eigensheaf support & Imaginary root multiplicity \\
Discriminant $\Delta \subset \cA_H$ & Singular opers & Wall-crossing locus \\
Langlands dual $G^\vee$ & Dual group & Koszul dual root datum \\
$T$-duality on fibres & Fourier--Mukai / Langlands & $\mathcal{G}(C,G)^! = \mathcal{G}(C,G^\vee)$ \\
6d $(2,0)$ on $C$ & Class $\mathcal{S}$ theory & Origin of $\mathcal{R}(C,G)$ \\
$4d \to 3d$ on $S^1$ & 3d mirror symmetry & Koszul duality at level of targets \\
$3d \to 2d$ on $S^1$ & 2d sigma model & Character of $\mathcal{G}(C,G)$ \\
KW twist at $t=1$ & Geometric Langlands twist & $\Etwo$-enhancement \\
Coulomb branch & Hitchin base & Cartan of $\mathcal{G}(C,G)$ \\
BPS states & Hecke operators & Root vectors of $\mathcal{G}(C,G)$ \\
Wall-crossing (KS) & Stokes data & Gauge equivalence of MC elements \\
FF center at $k = -h^\vee$ & Opers & Center of $\mathcal{G}(C,G)_{\mathrm{crit}}$ \\
\hline
\end{tabular}
\end{center}


% ==========================================================================
\section{Open problems}
\label{sec:open-problems}
% ==========================================================================

\begin{enumerate}[label=\textbf{OP\arabic*.}]
  \item \textbf{Explicit BKM for $G \neq \mathrm{SL}_2$.}  Construct the generalised BKM superalgebra $\frakg_{\mathcal{R}(C,G)}$ explicitly for $G = \mathrm{SL}_3$ and $g = 2$.  Identify the denominator identity as an automorphic form.  The Hitchin base is $H^0(C, K_C^2) \oplus H^0(C, K_C^3) \cong \mathbb{C}^3 \oplus \mathbb{C}^5 = \mathbb{C}^8$, and $\dim_\mathbb{C} \cM_H = 16$.

  \item \textbf{Wall-crossing as gauge equivalence.}  Make precise the identification of the KS wall-crossing formula with gauge equivalence of MC elements in the shadow obstruction tower $\Theta_\cA^{\leq r}$.

  \item \textbf{The automorphic form for general $(C, G)$.}  For $K3 \times E$, the automorphic form is $\Delta_5$.  What is the explicit automorphic form for $\mathcal{G}(C, G)$ when $G$ is not $\mathrm{SL}_2$ or when $g > 2$?  Is there a ``generalised Igusa cusp form'' for each pair $(g, G)$?

  \item \textbf{Non-simply-laced groups and outer automorphisms.}  The 6d $(2,0)$ theory is defined only for simply-laced $\frakg$.  For non-simply-laced $G$ (types $B_n, C_n, F_4, G_2$), the class $\mathcal{S}$ construction uses outer automorphism twists.  Formulate the quantum vertex chiral group $\mathcal{G}(C, G)$ in these cases and verify that Koszul duality still exchanges $G$ and $G^\vee$.

  \item \textbf{Quantum geometric Langlands from $\Etwo$-Koszul duality.}  Derive the Gaitsgory QGL equivalence $\Rep^{\Etwo}_\kappa(\mathcal{G}(C, G)) \simeq \Rep^{\Etwo}_{\kappa^\vee}(\mathcal{G}(C, G^\vee))$ from the $\Etwo$-bar-cobar adjunction of Volume~III.

  \item \textbf{Higher genus and the genus tower.}  The genus-$h$ partition function $Z(\Sigma_h; \cM_H(C, G))$ should be a genus-$h$ amplitude of $\mathcal{G}(C, G)$.  The genus tower (genus 0 = tree level, genus 1 = one-loop, genus $h$ = shadow depth $h$) should match the shadow obstruction tower of Volume~I.  Verify this at genus $h = 2$ for $G = \mathrm{SL}_2$, $g = 2$, where the genus-2 amplitude should involve $\Delta_5$ on $\mathbb{H}_2$.

  \item \textbf{Relation to the toric CY3 tower.}  The toric CY3 examples (the toric CoHA chapter of Vol.~III) give quantum vertex chiral groups from affine super Yangians.  Is there a ``deformation'' connecting the Hitchin quantum vertex chiral group $\mathcal{G}(C, G)$ to the toric one $\mathcal{G}(X_\Sigma)$?  Physically, this should correspond to a geometric transition (conifold transition or flop) relating the Hitchin CY geometry to a toric CY3.

  \item \textbf{Arithmetic Langlands.}  The geometric Langlands programme has arithmetic counterparts (the classical Langlands programme over number fields).  Does the quantum vertex chiral group framework extend to arithmetic settings?  In particular: for $C$ defined over $\mathbb{Q}$ (or a number field), is $\mathcal{G}(C, G)$ an object in ``arithmetic chiral algebra'' (in the sense of Beilinson--Drinfeld over $\Spec(\mathbb{Z})$)?
\end{enumerate}


% ==========================================================================
\section*{References}
% ==========================================================================

\begin{enumerate}[label={[\arabic*]}]
  \item N.~Hitchin, \textit{The self-duality equations on a Riemann surface}, Proc.\ London Math.\ Soc.\ \textbf{55} (1987), 59--126.
  \item N.~Hitchin, \textit{Stable bundles and integrable systems}, Duke Math.\ J.\ \textbf{54} (1987), 91--114.
  \item A.~Beilinson and V.~Drinfeld, \textit{Quantization of Hitchin's integrable system and Hecke eigensheaves}, preprint (1991).
  \item T.~Hausel and M.~Thaddeus, \textit{Mirror symmetry, Langlands duality, and the Hitchin system}, Invent.\ Math.\ \textbf{153} (2003), 197--229.
  \item R.~Donagi and T.~Pantev, \textit{Langlands duality for Hitchin systems}, Invent.\ Math.\ \textbf{189} (2012), 653--735.
  \item A.~Kapustin and E.~Witten, \textit{Electric-magnetic duality and the geometric Langlands program}, Commun.\ Number Theory Phys.\ \textbf{1} (2007), 1--236.
  \item D.~Gaiotto, \textit{$\cN = 2$ dualities}, JHEP \textbf{08} (2012), 034.
  \item B.~Feigin and E.~Frenkel, \textit{Affine Kac--Moody algebras at the critical level and Gelfand--Dikii algebras}, Int.\ J.\ Mod.\ Phys.\ A \textbf{7} Suppl.\ 1A (1992), 197--215.
  \item D.~Gaitsgory, \textit{Twisted Whittaker model and factorizable sheaves}, Selecta Math.\ \textbf{13} (2008), 617--659.
  \item D.~Gaitsgory et al., \textit{Proof of the geometric Langlands conjecture}, preprint (2024).
  \item V.~Gritsenko and V.~Nikulin, \textit{Siegel automorphic form corrections of some Lorentzian Kac--Moody Lie algebras}, Amer.\ J.\ Math.\ \textbf{119} (1997), 181--224.
  \item M.~Kontsevich and Y.~Soibelman, \textit{Stability structures, motivic Donaldson--Thomas invariants and cluster transformations}, preprint arXiv:0811.2435 (2008).
  \item M.~Rapcak, Y.~Soibelman, Y.~Yang, and G.~Zhao, \textit{Cohomological Hall algebras, vertex algebras, and instantons}, Commun.\ Math.\ Phys.\ \textbf{376} (2020), 1803--1873.
\end{enumerate}

\end{document}
